\subsection{Motivation}
\begin{prop}{Cauchy Criteria for Convergence}{}
	The following statements satisfy (1) $\iff$ (2) and (3) $\implies$ (2):
	\begin{enumerate}
		\item $(x_n)$ converges to some real $x \coloneqq \lim x_n$.
		\item $(x_n)$ is Cauchy convergent, that is, $\lim_{m,n \to \infty} |x_m - x_n| = 0$.
		\item $\sum_n \I_{\{|x_{n+1} - x_n| > \varepsilon_n\}} < \infty$ for any $(\varepsilon_n)$ such that $c \coloneqq \sum_n \varepsilon_n < \infty$.
	\end{enumerate}
\end{prop}

\subsection{Almost Sure Convergence}

\begin{defn}{Almost Sure Convergence}{}
	The sequence of random variables $(X_n)$ converges \textit{almost surely} ($\text{a.s.}$) to a random variable $X$ if
	\begin{equation*}
		\Prob(\lim X_n = X) = 1.
	\end{equation*}
	We denote the almost sure convergence by $X_n \xrightarrow{\text{a.s.}} X$.
\end{defn}

\begin{thm}{Almost Sure Convergence}{}
	The sequence of random variables $(X_n)$ converges almost surely to a random variable $X$ if and only if for every $\varepsilon > 0$
	\begin{equation}
		\lim_{m \to \infty} \Prob \left( \sup_{n \ge m} |X_n - X| > \varepsilon \right) = 0.
		\label{eq:AlmostSureConvergenceCriteria}
	\end{equation}
\end{thm}

\begin{prop}{Cauchy Criteria for Almost Sure Convergence}{}
	The following statements are equivalent:
	\begin{enumerate}
		\item $X_n \xrightarrow{\text{a.s.}} X$.
		\item $\lim_{n \to \infty} \Prob(\sup_{k \ge 0} |X_{n+k} - X_n| > \varepsilon) = 0$ for any choice of $\varepsilon > 0$.
		\item $\lim_{k \to \infty} \Prob(\sup_{m, n \ge k} |X_m - X_n| > \varepsilon) = 0$ for any choice of $\varepsilon > 0$.
		\item $\Prob(\lim_{m, n \to \infty} |X_m - X_n| = 0) = 1$; that is, $(X_n)$ is a.s. a Cauchy sequence.
	\end{enumerate}
\end{prop}

\begin{prop}{Markov's and Chebyshev's Inequality}{}
	For any random variable $Y$ and $\varepsilon > 0$, Markov's inequality holds:
	\begin{equation}
		\Prob(|Y| > \varepsilon) \le \frac{\E|Y|}{\varepsilon}.
		\label{eq:MarkovsInequality}
	\end{equation}
	As a consequence, we have Chebyshev's inequality, where $\mu \coloneqq \E X$:
	\begin{equation}
		\Prob(|X - \mu| > \varepsilon) \le \frac{\Var X}{\varepsilon^2}.
		\label{eq:ChebyshevsInequality}
	\end{equation}
\end{prop}

\begin{thm}{Kolmogorov's Inequality}{}
	Let $S_n \coloneqq \sum_{i=1}^n X_i$, where $X_1, X_2, \dots$ are independent with $0$ mean. Then, for every $\varepsilon > 0$:
	\begin{equation}
		\Prob \left( \sup_{k \le n} |S_k| > \varepsilon \right) \le \frac{\Var S_n}{\varepsilon^2}.
		\label{eq:KolmogorovsInequality}
	\end{equation}
\end{thm}

\begin{lem}{Borel--Cantelli}{}
	Let $(H_n)$ be a sequence of events. Then, the following hold:
	\begin{enumerate}
		\item $\sum_n \Prob(H_n) < \infty$ implies that $\Prob(\sum_n \I_{H_n} < \infty) = 1$.
		\item $\sum_n \Prob(H_n) = \infty$ implies that $\Prob(\sum_n \I_{H_n} < \infty) = 0$, provided that the events $(H_n)$ are pairwise independent.
	\end{enumerate}
\end{lem}

\subsection{Convergence in Probability}

\begin{defn}{Convergence in Probability}{}
	The sequence of random variables $(X_n)$ converges in probability to a random variable $X$ if, for all $\varepsilon > 0$,
	\begin{equation*}
		\lim_{n \to \infty} \Prob(|X_n - X| > \varepsilon) = 0.
	\end{equation*}
	We denote the convergence in probability by $X_n \xrightarrow{\Prob} X$.
\end{defn}

\begin{prop}{Convergence of a Subsequence}{}
	If $(X_n)$ converges to $X$ in probability, there is a subsequence $(X_{n_k})$ that converges to $X$ almost surely.
\end{prop}

\begin{prop}{Cauchy Criterion for Convergence in Probability}{}
	The sequence of random variables $(X_n)$ converges in probability if and only if for every $\varepsilon > 0$,
	\begin{equation}
		\lim_{m,n \to \infty} \Prob(|X_m - X_n| > \varepsilon) = 0.
		\label{eq:CauchyCriterionProbability}
	\end{equation}
\end{prop}

\subsection{Convergence in Distribution}
\begin{defn}{Convergence in Distribution}{}
	The sequence of random variables $(X_n)$ is said to \textit{convege in distribution} to a random variable $X$ with cdf $F$ provided that:
	\begin{equation}
		\lim_{n \to \infty} \Prob(X_n \le x) = F(x) \quad \text{for all } x \text{ such that } \lim_{a \to x} F(a) = F(x).
		\label{eq:ConvergenceInDistributionDef}
	\end{equation}
	We denote the convergence in distribution by $X_n \xrightarrow{d} X$.
\end{defn}

\begin{thm}{Characteristic Function and Convergence in Distribution}{}
	Suppose that $\psi_{X_1}(\vecb{r}), \psi_{X_2}(\vecb{r}), \dots$ are the characteristic functions of the sequence of $d$-dimensional random vectors $\vecb{X}_1, \vecb{X}_2, \dots$ and $\psi_X(\vecb{r})$ is the characteristic function of $\vecb{X}$. Then, the following three statements are equivalent:
	\begin{enumerate}
		\item $\lim_{n} \psi_{X_n}(\vecb{r}) = \psi_X(\vecb{r})$ for all $\vecb{r} \in \R^d$.
		\item $\vecb{X}_n \xrightarrow{d} \vecb{X}$.
		\item $\lim_n \E h(\vecb{X}_n) = \E h(\vecb{X})$ for all bounded continuous functions $h: \R^d \to \R$.
	\end{enumerate}
\end{thm}

\subsection{Convergence in $L^p$ Norm}

\begin{defn}{Convergence in $L^p$ Norm}{}
	The sequence of random variables $(X_n)$ converges in $L^p$ norm (for some $p \in [1, \infty]$) to a random variable $X$, if $\|X_n\|_p < \infty$ for all $n$, $\|X\|_p < \infty$, and
	\begin{equation*}
		\lim_{n \to \infty} \|X_n - X\|_p = 0.
	\end{equation*}
	We denote the convergence in $L^p$ norm by $X_n \xrightarrow{L^p} X$.
\end{defn}

\begin{prop}{Cauchy Criterion for Convergence in $L^p$ Norm}{CauchyCriteriaLp}
	The sequence of random variables $(X_n)$ converges in $L^p$ norm for $p \in [1, \infty]$ if and only if $\|X_n\|_p < \infty$ for all $n$ and
	\begin{equation*}
		\lim_{m,n \to \infty} \|X_m - X_n\|_p = 0.
	\end{equation*}
\end{prop}

\begin{prop}{Integrable Random Variable}{}
	A real-valued random variable $X$ is integrable if and only if
	\begin{equation}
		\lim_{b \to \infty} \E |X| \I_{\{|X| > b\}} = 0.
		\label{eq:IntegrableRandomVariableCondition}
	\end{equation}
\end{prop}

\begin{defn}{Uniform Integrability}{}
	A collection $\mathcal{K}$ of random variables is said to be \textit{uniformly integrable} ($\text{UI}$) if
	\begin{equation*}
		\lim_{b \to \infty} \sup_{X \in \mathcal{K}} \E |X| \I_{\{|X| > b\}} = 0.
	\end{equation*}
\end{defn}

\begin{prop}{Conditions for Uniform Integrability}{}
	Let $\mathcal{K}$ be a collection of random variables.
	\begin{enumerate}
		\item If $\mathcal{K}$ is a finite collection of integrable random variables, then it is UI.
		\item If $|X| \le Y$ for all $X \in \mathcal{K}$ and some integrable $Y$, then $\mathcal{K}$ is UI.
		\item $\mathcal{K}$ is UI if and only if there is an increasing convex function $f$ such that
		      \begin{equation*}
			      \lim_{x \to \infty} \frac{f(x)}{x} = \infty \quad \text{and} \quad \sup_{X \in \mathcal{K}} \E f(|X|) < \infty.
		      \end{equation*}
		\item $\mathcal{K}$ is UI if $\sup_{X \in \mathcal{K}} \E |X|^{1+\varepsilon} < \infty$ for some $\varepsilon > 0$.
		\item $\mathcal{K}$ is UI if and only if
		      \begin{enumerate}
			      \item $\sup_{X \in \mathcal{K}} \E |X| < \infty$, and
			      \item for every $\varepsilon > 0$ there is a $\delta > 0$ such that for every event $H$:
			            \begin{equation}
				            \Prob(H) \le \delta \implies \sup_{X \in \mathcal{K}} \E |X| \I_{H} \le \varepsilon.
				            \label{eq:UniformIntegrabilityEquivalence}
			            \end{equation}
		      \end{enumerate}
	\end{enumerate}
\end{prop}

\begin{thm}{$L^1$ Convergence and Uniform Integrability}{}
	A sequence $(X_n)$ of real-valued integrable random variables converges in $L^1$ if and only if it converges in probability and is uniformly integrable.
\end{thm}

\subsection{Law of Large Numbers and Central Limit Theorem}

\begin{thm}{Weak Law of Large Numbers}{WeakLawLargeNumbers}
	If $X_1, X_2, \dots$ are $\iid$ with finite expectation $\mu$, then for all $\varepsilon > 0$
	\begin{equation*}
		\lim_{n \to \infty} \Prob(|\xbar{X}_n - \mu| > \varepsilon) = 0.
	\end{equation*}
	In other words, $\xbar{X}_n \xrightarrow{\Prob} \mu$.
\end{thm}

\begin{thm}{Strong Law of Large Numbers}{StrongLawLargeNumbers}
	If $X_1, X_2, \dots$ are $\iid$ with expectation $\mu$ and $\E X^2 < \infty$, then for all $\varepsilon > 0$
	\begin{equation*}
		\lim_{m \to \infty} \Prob \left( \sup_{n \ge m} |\xbar{X}_n - \mu| > \varepsilon \right) = 0.
	\end{equation*}
	In other words, $\xbar{X}_n \xrightarrow{\text{a.s.}} \mu$.
\end{thm}

\begin{thm}{Central Limit Theorem}{CentralLimitTheorem}
	If $X_1, X_2, \dots$ are $\iid$ with finite expectation $\mu$ and finite variance $\sigma^2$, then for all $x \in \R$,
	\begin{equation*}
		\lim_{n \to \infty} \Prob \left( \frac{\xbar{X}_n - \mu}{\sigma/\sqrt{n}} \le x \right) = \Phi(x),
	\end{equation*}
	where $\Phi$ is the cdf of the standard normal distribution. In other words, $\sqrt{n}(\xbar{X}_n - \mu)/\sigma \xrightarrow{d} \mathcal{N}(0, 1)$.
\end{thm}
